\scanpage{133}
% archive foliation: 132
\deleted{\unsure{de sous-groupe}} isomorphe à $\mathbb U=\{z\in\mathbb C\mid |z|=1\}$, \inserted{tel que l'action sur $V_{\mathbb C}$ ait exactement $2$ sous-espaces propres}
\margin{(c'est automatique si $\operatorname{rg}_{\mathbb R}V=2$)}

$2^\circ)$ D'un des deux éléments d'ordre $4$ de $U$, ou ce qui revient au $=$, d'une \uline{orientation} de $U$, ou encore \inserted{(si $\operatorname{rg}_{\mathbb R}V=2$, plus gén. si $\operatorname{rg}_{\mathbb R}V=2n$, $n$ impair)} d'une orientation de $V$.
\margin{ou encore d'un \inserted{des deux} isomorphismes $\mathbb U\xrightarrow{\sim}U$, $U\xrightarrow{\sim}\mathbb U$}

--- l'élément d'ordre $4$ de $U$ associé à une orientation $\omega$ de \inserted{$U$ (ou de} $V$\inserted{)} sera noté $u_\omega$, c'est donc la rotation de $+\pi/2$ pour l'orientation \inserted{$\omega$} de $V$ et toute structure \deleted{\unsure{euclidienne}}\inserted{euclidienne} sur $V$ choisie invariante sous $U$ (\uline{NB} \inserted{si $\operatorname{rg}_{\mathbb R}V=2$,} une telle \deleted{\unsure{structure}}\inserted{métrique} est unique à scalaire multiplicatif près).

Alors $V_{\mathbb C}$ sous l'action de $U$ se décompose en deux sous-espaces propres
\begin{equation}
V_{\mathbb C}=\bigoplus_{\omega\in\underline\omega}L_\omega
\tag{\unsure{14|16}}
\end{equation}
où pour $\omega\in\underline\omega(U)=\underline\omega$, $L_\omega$ est caractérisé par le fait que l'isomorphisme
\begin{equation}
\varphi_\omega:U\xrightarrow{\sim}\mathbb U
\qquad
g(z)=\varphi_\omega(g)\cdot z\quad\text{pour }z\in L_\omega
\tag{15}
\end{equation}
défini par l'action de $U$ sur $L_\omega$, \unsure{corresponde|correspond} à l'orientation $\omega$. La donnée d'une structure complexe sur $V$ \inserted{compatible avec $U$} équivaut alors à celle d'un des deux isomorphismes $\varphi_\omega$ (permettant de faire opérer $\mathbb U$ sur $V$, donc aussi $\mathbb R^{\times+}\times\mathbb U\simeq\mathbb C^\times$, \deleted{\unsure{donc}} $\mathbb C\ldots$). Ceci dit, pour $\deleted{\unsure{z\in\mathbb C}}\ z\in L_\omega$, on aura
\[
\operatorname{Re}(\varphi_\omega(g)z)=\operatorname{Re}(gz)=g\operatorname{Re}(z)
=\text{produit de }\operatorname{Re}(z)\text{ par }\varphi_\omega(g)
\text{ pour la structure complexe de }V\text{ associée à }\varphi_\omega.
\]
